{
{\Large {\bfseries П}ідсилення {\bfseries с}вітла {\bfseries в}имушеним {\bfseries в}ипромінюванням ({\bfseries LASER})}
	\rput(0.3in,0.04in){
	\psset{	linecolor=yellow,
		fillcolor=yellow,
		linewidth=0.6pt}
		\newlength{\nLASERsize}	\setlength{\nLASERsize}{0.12in}
		\newlength{\nLASERray}	\setlength{\nLASERray}{\nLASERsize}
		\newlength{\nLASERsray}	\setlength{\nLASERsray}{0.75\nLASERsize}
		\multido{\rLASER=0+30}{6}{%Large rays
			\rput{\rLASER}(0,0){\psline(-\nLASERray,0in)(\nLASERray,0in)}
		}
		\multido{\rLASER=15+30}{6}{%Small rays
			\rput{\rLASER}(0,0){\psline(-\nLASERsray,0in)(\nLASERsray,0in)}
		}
		\pscircle(0,0){0.4\nLASERsize} %Inner circle
		\rput(0,0){\psline[linewidth=0.7pt](0in,0in)(2\nLASERray,0in)}%Long line
	}
\begin{itemize}

% \item LASER is an acronym for {\bfseries L}ight {\bfseries A}mplification by {\bfseries S}timulated {\bfseries E}mission of {\bfseries R}adiation.
\item LASER --- це пристрій, що генерує когерентне монохроматичне EMR високої інтенсивності.
\item За допомогою відповідного обладнання будь-яке EMR, таке як видиме світло, ультрафіолет, рентгенівські та гамма-промені, може працювати як LASER. MASER функціонує подібно до LASER і працює з мікрохвилями та радіохвилями.

%\item Junk/reword: Photons can be brought together and the resultant sum of forces is seen. If the polarization and phase are the same for the photons the the resultant light appears to be amplified. Some LASERs achieve this by constantly pumping photons into a filter to get the same frequency then bouncing those filtered photons off opposing mirrors to get the phase and polarization just right. Once the light has been ``amplified" enough it will escape from one of the mirrors.

\end{itemize}
}
